\documentclass[12pt,leqno]{article}
\usepackage{amsmath,amssymb,amsthm,colortbl,xcolor,nicefrac,setspace}
\usepackage[T1]{fontenc}
\usepackage[expert,seriftt,lucidasmallscale]{lucidabr}
\usepackage{mathrsfs,hyperref,color,graphicx,needspace}
\hypersetup{hidelinks,hypertexnames=false}
\setlength{\textwidth}{6.5in}
\setlength{\textheight}{9in}
\setlength{\oddsidemargin}{0mm}
\setlength{\evensidemargin}{0mm}
\setlength{\topmargin}{0mm}
\setlength{\headheight}{0mm}
\setlength{\headsep}{0mm}
\setlength{\topskip}{0mm}
\setlength{\hoffset}{0in}
\setlength{\voffset}{0in}
\newcommand{\AND}{\ensuremath{\mathbin{\&}}}
\newcommand{\OR}{\ensuremath{\vee}}
\newcommand{\IFF}{\ensuremath{\leftrightarrow}}
\newcommand{\IC}{\ensuremath{\rightarrow}}
\newcommand{\NOT}{\ensuremath{\text{$\thicksim$}}}
\newcommand{\Pure}{\mathrm{Pure}}
\newcommand{\Fun}{\mathrm{Fun}}
\newcommand{\funp}{\mathrm{fun}'}
\newcommand{\Stage}{\mathrm{Stage}}
\newcommand{\Pre}{\mathrm{Pre}}
\newcommand{\Clos}{\mathrm{Cl}}
\newcommand{\Th}{\mathcal T}
\newcommand{\cZ}{\mathcal Z}
\newcommand{\Kbot}{K_{\bot}}
\newcommand{\Ktop}{K_{\top}}
\newtheorem{theorem}{Theorem}
\newtheorem{lemma}{Lemma}
\renewcommand{\qedsymbol}{\(\blacksquare\)}
\onehalfspacing

\begin{document}
\begin{center}
{\Large Purity of Pure and Recombination}

\smallskip
11 November 2026
\end{center}

We prove that, over the background theory $T_0$, Purity of Pure at type
$t\to t$ is inconsistent with unary Recombination and the existence of
a fundamental proposition. We also prove that $T_0$ and this single instance of Purity of Pure rule
out a $\funp$ proposition.

We use Goodman's notation and put $u=t\to t$. Lower-case variables
$p,q,r,s$ range over propositions. Upper-case variables $V,W,X,Y,B$ range
over unary operators of type $u$. The symbol $\cZ$ ranges over properties
of operators, of type $u\to t$. We write
\[
 \Box A=(A=\top),\qquad
 V\subseteq W\quad\text{for}\quad\forall p\,(Vp\IC Wp).
\]
The inclusion $V\subseteq W$ says that every proposition satisfying $V$ satisfies $W$.
When we need necessary inclusion, we write $\Box(V\subseteq W)$.

The background theory $T_0$ includes Classicism, the Purity schema and
Application closure. The Purity schema says that every closed logical
term is pure. Application closure is the schema
\[
 \Pure_{\sigma\to\tau}(F)\AND\Pure_\sigma(a)
       \IC\Pure_\tau(Fa).
\]
Let $\Th$ be the axiom extension of $T_0$ by
\begin{equation}
 \Pure_{u\to t}(\Pure_u).
 \tag{PP$_u$}\label{pp}
\end{equation}
The Rule of Equivalence applies to theorems of this extension.
Consequently, if $\Th\vdash A$, then $\Th\vdash A=\top$, and thus
$\Th\vdash\Box A$. We use this rule of Necessitation, the S4 laws,
Boolean identities, $\beta$-conversion and typed universal instantiation.
In Lemma~\ref{bridge} we explicitly mark the application of Necessitation
to a theorem whose proof uses \eqref{pp}.

\begin{theorem}\label{main}
The theory $\Th$ proves that no proposition $r$ satisfies
\begin{equation}
 \forall X^u\bigl(\Pure(X)\IC
        (\Box(Xr)\IC\forall p\,Xp)\bigr).
 \tag{R}\label{R}
\end{equation}
Hence $\Th\vdash\NOT\exists r\,\funp(r)$, and
\[
 \Th+\text{unary Recombination}\vdash\NOT\exists r\,\Fun(r).
\]
\end{theorem}

\Needspace{6\baselineskip}
Unary Recombination requires \eqref{R} of every fundamental proposition.
The definition of $\funp(r)$ is
\[
 \forall X^uY^u\bigl((\Pure(X)\AND\Pure(Y))
           \IC(Xr=Yr\IC X=Y)\bigr).
\]
We prove the theorem by constructing two unary operators, $L$ and $P$.
The closed logical operator $L$ has the modal properties needed for the
diagonal argument. The operator $P$ collects the instances of a
definable class of pure operators called stages. PP makes $P$ pure and
allows us to prove $L\subseteq P$. For each proposed $r$, we then obtain
a proposition $q=r\AND Lr$ and a pure stage $Y_0$ that omits $q$ but
whose successor contains it. This produces a pure operator $X$ with
$\Box(Xr)$ for which $X(\NOT q)$ is false, contradicting \eqref{R}.

\medskip
\noindent\textbf{A logical fixed-point construction.}
Define
\begin{align}
 \Sigma&=\lambda Y^u.\lambda p.
  \bigl(\Box(p\OR Yp\OR Y(\NOT p))
       \OR\Box(\NOT p\OR Yp\OR Y(\NOT p))\bigr),\label{sigma}\\
 \Pre(W)&=\Box(\Sigma W\subseteq W),\label{pre}\\
 L&=\lambda p.\forall W^u\,(\Pre(W)\IC Wp).\label{L}
\end{align}
Here $\Sigma$ has type $u\to u$, $\Pre$ has type $u\to t$, and
$L$ has type $u$. Each abbreviation expands to a closed logical term.
The condition $\Pre(W)$ says that $W$ necessarily contains its
$\Sigma$-successor. The operator $L$ applies to precisely the
propositions satisfying every such $W$.

\Needspace{8\baselineskip}
\begin{lemma}\label{sigma-properties}
The following are logical theorems.
\begin{align}
 (\Sigma Y)(\NOT p)&\IFF(\Sigma Y)p,\label{sigma-sym}\\
 (\Sigma Y)p&\IC\Box((\Sigma Y)p),\label{sigma-pers}\\
 \Box(V\subseteq W)&\IC\Box(\Sigma V\subseteq\Sigma W).
 \label{sigma-mono}
\end{align}
\end{lemma}
\begin{proof}
Replacing $p$ by $\NOT p$ in \eqref{sigma} interchanges the two boxed
disjuncts and interchanges $Yp$ with $Y(\NOT p)$ inside each box.
Double negation and commutativity of disjunction give
\eqref{sigma-sym}. For \eqref{sigma-pers}, S4 proves
\[
 \Box A\OR\Box B\IC\Box(\Box A\OR\Box B).
\]
Apply this to the two disjuncts in \eqref{sigma}.

For monotonicity, suppose that $\Box(V\subseteq W)$. Boxed universal
instantiation gives both $\Box(Vp\IC Wp)$ and
$\Box(V(\NOT p)\IC W(\NOT p))$. The normal modal rules therefore
take each disjunct of $(\Sigma V)p$ to its corresponding disjunct in
$(\Sigma W)p$. Generalizing $p$ and discharging the assumption proves
\begin{equation}
 \Box(V\subseteq W)\IC(\Sigma V\subseteq\Sigma W).
 \label{mono-local}
\end{equation}
Necessitate the universal closure of this theorem and instantiate inside
the box at $V,W$. Axiom $4$ gives
$\Box(V\subseteq W)\IC\Box\Box(V\subseteq W)$, so $K$ gives
the boxed conclusion in \eqref{sigma-mono}.
\end{proof}

\begin{lemma}\label{fixed}
Logic alone proves
\begin{align}
 Lp&\IFF(\Sigma L)p,\label{fixed-law}\\
 Lp&\IC\Box(Lp),\label{L-pers}\\
 L(\NOT p)&\IFF Lp.\label{L-sym}
\end{align}
It also proves $\Pre(W)\IC\Box(L\subseteq W)$ for every $W$.
\end{lemma}
\begin{proof}
Instantiating the quantifier in \eqref{L} gives the logical theorem
$\Pre(W)\IC L\subseteq W$. Its necessitation, together with
$\Pre(W)\IC\Box\Pre(W)$ from axiom $4$, gives
\begin{equation}
 \Pre(W)\IC\Box(L\subseteq W).
 \label{least}
\end{equation}
We first show $\Sigma L\subseteq L$. Suppose that $(\Sigma L)p$ and
let $W$ satisfy $\Pre(W)$. By \eqref{least} and
\eqref{mono-local}, $(\Sigma W)p$. Factivity applied to
$\Pre(W)$ then gives $Wp$. Discharging $\Pre(W)$ and generalizing
$W$ gives $Lp$ by \eqref{L}. After discharging $(\Sigma L)p$ and generalizing $p$,
we have proved $\Sigma L\subseteq L$ as a logical theorem.
Necessitation therefore gives
\[
 \Box(\Sigma L\subseteq L),
\]
which is $\Pre(L)$.

Now apply the stronger monotonicity law \eqref{sigma-mono} with
$V=\Sigma L$ and $W=L$. It gives
\[
 \Box(\Sigma(\Sigma L)\subseteq\Sigma L),
\]
which is $\Pre(\Sigma L)$. Instantiating \eqref{L} at
$W=\Sigma L$ thus proves $L\subseteq\Sigma L$.
Together the two inclusions give \eqref{fixed-law}.

The fixed-point biconditional is a logical theorem, so its necessitation
is available. From $Lp$, it gives $(\Sigma L)p$, then
$\Box((\Sigma L)p)$ by \eqref{sigma-pers}, and finally $\Box(Lp)$
by $K$ and the necessitated reverse implication. This proves
\eqref{L-pers}. The symmetry law \eqref{L-sym} follows by applying
\eqref{fixed-law} on both sides of \eqref{sigma-sym} with $Y=L$.
\end{proof}

\Needspace{8\baselineskip}
\begin{lemma}\label{diagonal}
Put $g=\lambda p.(p\AND Lp)$. For each proposition $r$, put $q=gr$.
Then logic proves $Lq$ and $g(\NOT q)=\NOT q$.
\end{lemma}
\begin{proof}
Conjoining $Lr$ to $r$ gives a proposition $q$ for which we can first
prove $q\IC Lq$. Assume $q$, so $Lr$. By \eqref{L-pers},
$\Box(Lr)$, which by definition is $Lr=\top$. Hence
\[
 q=r\AND Lr=r\AND\top=r.
\]
Substitution in $Lr$ gives $Lq$.
Discharging $q$ and generalizing $r$ proves
\[
 \forall r\bigl((r\AND Lr)\IC L(r\AND Lr)\bigr).
\]
Necessitation and boxed universal instantiation therefore give
$\Box(\NOT q\OR Lq)$. By weakening inside the box we obtain
\[
 \Box(\NOT q\OR Lq\OR L(\NOT q)).
\]
This is the second disjunct of $(\Sigma L)q$, so
\eqref{fixed-law} gives $Lq$. Symmetry gives $L(\NOT q)$, and
\eqref{L-pers} gives $L(\NOT q)=\top$. Thus
\[
 g(\NOT q)=\NOT q\AND L(\NOT q)
          =\NOT q\AND\top=\NOT q.\qedhere
\]
\end{proof}

The identities $gr=q$ and $g(\NOT q)=\NOT q$ will determine the
values of our final operators at the two arguments used in the
contradiction. Both $L$ and $g$ are closed logical operators and
are therefore pure by the Purity schema.

\medskip
\noindent\textbf{The pure stages.}
Let $\Kbot=\lambda p.\bot$. A property $\cZ$ of operators is called
\emph{closed} when it satisfies the following conditions.
\begin{enumerate}
\item $\cZ(\Kbot)$.
\item For every pure $Y$, if $\cZ(Y)$, then $\cZ(\Sigma Y)$.
\item For every pure $B$, if each proposition $s$ with $Bs$ has a
pure operator $Y$ satisfying $\cZ(Y)$, $Y\subseteq B$, and $Ys$,
then $\cZ(B)$.
\end{enumerate}
The third condition admits a pure $B$ whenever its extension is covered
by the extensions of pure members of $\cZ$ that it contains. We call
this the union condition. Its formal expression is
\[
 \forall B\Bigl(\Bigl(\Pure(B)\AND\forall s
  \bigl(Bs\IC\exists Y(\Pure(Y)\AND\cZ(Y)\AND Y\subseteq B\AND Ys)\bigr)
  \Bigr)\IC\cZ(B)\Bigr).
\]
Let $\Clos(\cZ)$ be the conjunction of the three conditions, and define
\begin{align}
 \Stage(Y)&=\Pure(Y)\AND\forall\cZ\,(\Clos(\cZ)\IC\cZ(Y)),
 \label{stage}\\
 P&=\lambda p.\exists Y\,(\Stage(Y)\AND Yp).
 \label{P}
\end{align}
Here $\Clos$ has type $(u\to t)\to t$, $\Stage$ has type $u\to t$,
and $P$ has type $u$. The quantifier over $\cZ$ ranges over properties
of type $u\to t$. The first conjunct in \eqref{stage} makes every
stage pure. Equation~\eqref{P} gives $Y\subseteq P$ for every stage $Y$.

\Needspace{12\baselineskip}
\begin{lemma}\label{induct}
The property $\Stage$ is closed. If a property $\Psi$ of type $u\to t$
satisfies the following three conditions, then every stage satisfies $\Psi$.
\begin{enumerate}
\item $\Psi(\Kbot)$.
\item If $Y$ is a stage satisfying $\Psi$, then $\Psi(\Sigma Y)$.
\item If $B$ is pure and each $s$ with $Bs$ has a stage $Y\subseteq B$
with $\Psi(Y)$ and $Ys$, then $\Psi(B)$.
\end{enumerate}
Parameters are allowed in $\Psi$. Bound variables are chosen fresh for them.
\end{lemma}
\begin{proof}
The operator $\Kbot$ is pure by the Purity schema and satisfies every
closed property by its first condition. Hence it is a stage.
If $Y$ is a stage, then $\Sigma Y$ is pure by logical purity of
$\Sigma$ and Application closure. Every closed property contains $Y$
and hence $\Sigma Y$, making $\Sigma Y$ a stage.

For the union condition, suppose that $B$ is pure and every $s$ with
$Bs$ has a stage $Y\subseteq B$ with $Ys$. Fix a closed property
$\cZ$. Each such $Y$ is pure and satisfies $\cZ$, by
\eqref{stage}. The third condition for $\cZ$ gives $\cZ(B)$.
Generalizing $\cZ$ proves that $B$ is a stage. Thus $\Stage$ is closed.

Now suppose that $\Psi$ satisfies the displayed conditions and put
\[
 \mathcal W=\lambda Y^u.(\Stage(Y)\AND\Psi(Y)).
\]
The first and second closure conditions for $\mathcal W$ follow from
those for $\Stage$ and the corresponding conditions on $\Psi$.
For the third, the witnesses in $\mathcal W$ are stages satisfying
$\Psi$, so the union condition for $\Stage$ gives $\Stage(B)$
and the third condition on $\Psi$ gives $\Psi(B)$.
Thus $\Clos(\mathcal W)$. Instantiating \eqref{stage} at this
typed $\lambda$-term shows that every stage satisfies $\Psi$.
\end{proof}

\begin{lemma}\label{stage-properties}
Every stage $Y$ satisfies
\begin{equation}
 \forall p\,(Y(\NOT p)\IFF Yp).
 \label{stage-sym}
\end{equation}
If a stage $Y$ satisfies $Yp$, then there is a stage $Y_0$ such that
\begin{equation}
 \NOT Y_0p\AND(\Sigma Y_0)p.
 \label{predecessor}
\end{equation}
\end{lemma}
\begin{proof}
For symmetry, apply Lemma~\ref{induct} to the property
$\lambda Y.\forall p(Y(\NOT p)\IFF Yp)$. The base case follows
from $\Kbot p=\bot$. Every successor is symmetric by
\eqref{sigma-sym}. For the union condition, if $Bs$, choose a
symmetric stage $Y\subseteq B$ with $Ys$. Then $Y(\NOT s)$,
so $B(\NOT s)$. Applying the same argument to $\NOT s$ and using
double negation gives the converse.

For the second assertion fix $p$ and apply Lemma~\ref{induct} to
\[
 \Psi=\lambda Y^u.\Bigl(Yp\IC\exists Y_0^u
       (\Stage(Y_0)\AND\NOT Y_0p\AND(\Sigma Y_0)p)\Bigr).
\]
The base case is vacuous. In the successor case suppose that
$(\Sigma Y)p$. If $Yp$, the induction hypothesis supplies $Y_0$.
If $\NOT Yp$, take $Y_0=Y$. For the union condition, an instance
$Bp$ has a stage $Y\subseteq B$ with $Yp$ and $\Psi(Y)$.
The latter supplies the desired $Y_0$. Generalizing $p$ proves
the assertion.
\end{proof}

\begin{lemma}\label{bridge}
The theory $\Th$ proves that $P$ is pure, $P$ is a stage, and
$L\subseteq P$.
\end{lemma}
\begin{proof}
The term defining $P$ uses the single nonlogical constant $\Pure_u$.
Let $c$ have type $u\to t$, and let $\Clos_c$ be the closure formula
with $c$ replacing every occurrence of $\Pure_u$. The term
\[
 M=\lambda c^{u\to t}.\lambda p.
    \exists Y^u\Bigl(cY\AND
       \forall\cZ^{u\to t}(\Clos_c(\cZ)\IC\cZ(Y))\AND Yp\Bigr)
\]
is closed and logical, of type $(u\to t)\to u$.
By $\beta$-conversion, $M(\Pure_u)=P$.
Logical purity gives $\Pure(M)$, while \eqref{pp} gives
$\Pure_{u\to t}(\Pure_u)$. Application closure at input type
$u\to t$ and output type $u$ therefore gives $\Pure(P)$.

For each $p$ with $Pp$, definition~\eqref{P} supplies a stage $Y$
with $Yp$, and $Y\subseteq P$ holds by the same definition.
Since $P$ is pure, the union condition for $\Stage$ applies to $P$
and proves $\Stage(P)$. Successor closure gives $\Stage(\Sigma P)$,
and hence
\begin{equation}
 \Sigma P\subseteq P.
 \label{P-pre}
\end{equation}
This is a closed theorem of the axiom extension $\Th$. Its proof has
used the Purity schema, Application closure, and \eqref{pp}, with all
temporary assumptions discharged. Necessitation therefore gives
\[
 \Box(\Sigma P\subseteq P).
\]
Thus $\Pre(P)$ holds, by \eqref{pre}.
Instantiating the quantifier in \eqref{L} at $W=P$ now gives
$L\subseteq P$.
\end{proof}

PP supplies the inclusion $L\subseteq P$ used in the final argument.
Lemma~\ref{diagonal} supplies $Lq$, and Lemma~\ref{bridge} gives $Pq$.
The definition of $P$ then gives a stage containing $q$. Lemma~\ref{stage-properties}
then supplies a pure $Y_0$ for which \eqref{predecessor} holds at $q$.

\Needspace{6\baselineskip}
\begin{proof}[Proof of Theorem~\ref{main}]
Suppose that $r$ satisfies \eqref{R}. Put $q=r\AND Lr=gr$.
By Lemmas~\ref{diagonal} and~\ref{bridge}, $Pq$.
Choose a stage containing $q$. Lemma~\ref{stage-properties} gives
a stage $Y_0$ with
\[
 \NOT Y_0q,\qquad(\Sigma Y_0)q.
\]
It is pure and symmetric, so $\NOT Y_0(\NOT q)$ also holds.
We consider the two disjuncts of $(\Sigma Y_0)q$.

\emph{Case 1.} Suppose that
\[
 \Box(q\OR Y_0q\OR Y_0(\NOT q)).
\]
Factivity gives the disjunction inside the box. Since its last two
disjuncts are false, $q$ holds. Define
\[
 X_+=\lambda p.\bigl(gp\OR Y_0(gp)\OR Y_0(\NOT(gp))\bigr).
\]
Its purity follows from the Purity schema and Application closure,
using the closed logical term
\[
 D_+=\lambda Y^u.\lambda p.
       \bigl(gp\OR Y(gp)\OR Y(\NOT(gp))\bigr)
\]
applied to the pure $Y_0$. Here $g$ expands to a closed logical term,
so $D_+$ does too, with type $u\to u$.
The identity $gr=q$ identifies $X_+r$ with the entire disjunction
in the box, giving $\Box(X_+r)$.
Condition \eqref{R} yields $\forall p\,X_+p$, whose instance at
$\NOT q$ is $X_+(\NOT q)$.
Using $g(\NOT q)=\NOT q$ and double negation, we obtain
\[
 \NOT q\OR Y_0(\NOT q)\OR Y_0q.
\]
Each disjunct is false, which is a contradiction.

\emph{Case 2.} Suppose that
\[
 \Box(\NOT q\OR Y_0q\OR Y_0(\NOT q)).
\]
Factivity and the same two local falsities give $\NOT q$. Define
\[
 X_-=\lambda p.\bigl(\NOT(gp)\OR Y_0(gp)\OR Y_0(\NOT(gp))\bigr).
\]
This is pure by logical purity of the corresponding term
$D_-=\lambda Y^u.\lambda p.
 (\NOT(gp)\OR Y(gp)\OR Y(\NOT(gp)))$ and Application closure.
Again $gr=q$ gives $\Box(X_-r)$. Condition \eqref{R} yields
$\forall p\,X_-p$, and therefore $X_-(\NOT q)$.
Using $g(\NOT q)=\NOT q$ and double negation, we obtain
\[
 q\OR Y_0(\NOT q)\OR Y_0q.
\]
All three disjuncts are false. This completes the contradiction
in both cases, and proves that no $r$ satisfies \eqref{R}.

\Needspace{8\baselineskip}
For the first consequence suppose that $\funp(r)$ and that $X$ is
pure with $\Box(Xr)$. The operator $\Ktop=\lambda p.\top$ is pure
by the Purity schema. Since
\[
 Xr=\top=\Ktop r,
\]
the definition of $\funp(r)$ yields $X=\Ktop$. It follows that
$Xp=\top$, hence $Xp$, for every $p$. Thus $r$ satisfies
\eqref{R}, contradicting the result just proved.

Finally, unary Recombination states
\[
 \forall Xr\bigl((\Pure(X)\AND\Fun(r))
              \IC(\Box(Xr)\IC\forall p\,Xp)\bigr).
\]
Every fundamental proposition would therefore satisfy \eqref{R}.
This proves the second consequence.
\end{proof}

\end{document}
